\section{The quadratic-character extension of the original-level exclusion}
\label{sec:cubic-quadratic-typeii-exclusion}

This is a source-scoped extension lemma for incorporation with AC-G116,
not a second copy of its all-residue certificate or measure argument.
It concerns an actual original point and its exact local conditions.
It makes no assertion about a seed with additional auxiliary tame type.

Fix $p=269$, $K=\mathbb Q(\sqrt{-3})$, and $\eta=\eta_K$.
Let $D$ be the literal original \texttt{Dat} for $e,N,\psi$,
with all \texttt{GD}, local, coefficient and modularity inputs retained
in AC-G68-point.  In particular the original $e_O=1$ condition is
retained when it is used to obtain modularity.  Assume
\[
                         \psi^2=1.
\]
This includes the trivial character.  No condition $p\nmid N$ is added.
The character $c_p$ is the actual action on roots of unity, and the
source \texttt{Newform.Att} evaluates geometric Frobenius.

\begin{lemma}[Exact local reduction for an arbitrary quadratic original character]
Suppose an actual original Type II point supplies the endpoint
$V:G_{\mathbb Q}\to\operatorname{GL}_2(E)$ of AC-G68-point, with its
original integral lattice, determinant $\eta$, and residual
semisimplification
\[
 \overline V^{\rm ss}
   =\eta\overline\psi\overline c_p^{\,e}
          \oplus\overline\psi^{-1}\overline c_p^{-e}.
\]
Let $f$ be its actual attached newform, so
$R_f=Vc_p^{-e}$ and $k=2e+1$.
The primitive automorphic twist
\[
        \pi_g=\pi_f\otimes((\psi\eta)\circ\det)
\]
has weight $k$, nebentypus $\eta$, and primitive conductor exactly $3$.
Its representation is
\[
 R_g=V\psi\eta c_p^{-e},\qquad
 \overline R_g^{\rm ss}=1\oplus\eta\overline c_p^{-2e}.
\]
It is ordinary and crystalline at $269$.  Its actual primitive
coefficients satisfy $a_3(g)\equiv a_{269}(g)\equiv1$ in the residue
field.  These conclusions also hold when $\psi$ has conductor divisible
by $269$.
\end{lemma}

\begin{proof}
\emph{The literal conductor-primary character.}
The complete definition of \texttt{GD} in \texttt{Counts.lean}
contains $\operatorname{comp}_p(\psi)^2=1$, not $p\nmid N$.
The CRT idempotent in
\texttt{DirichletPrimeToLPartMulLPowerPart.lean} is one on the
$p$-power factor of the modulus and zero on its complementary factor;
the ensuing homomorphism on units sends $(u_p,u^{(p)})$ to $(u_p,1)$.
Thus $\operatorname{comp}_p\psi$ is precisely the character induced
by the $p$-power factor.  The literal \texttt{RhoKApply.lean},
\texttt{Dat.loc}, sets
\[
 \tau=(\texttt{Chr.dirG}(\operatorname{comp}_p\psi))\circ D.\iota.
\]
It is not a newly chosen quadratic twist.  Its restriction to inertia
is $\psi|_{I_p}$, and $\psi/\tau$ is unramified at $p$.
Since $p$ is odd, a character of order dividing two is trivial on
$1+p\mathbb Z_p$.  Hence the primitive conductor of this component is
$1$ or $p$, even if the displayed modulus contains a higher power.

\emph{Remove the finite $p$-type before crystalline descent.}
The full proof of \texttt{admissible\_of\_typeII} constructs a stable
line on the rational decomposition group, not merely over $K_p$.
Write its two characters as $\delta_1,\delta_2$; then
\[
 \delta_1|_{I_p}=\psi c_p^e,\qquad
 \delta_2|_{I_p}=\psi^{-1}c_p^{-e},\qquad
 \delta_1\delta_2=\eta.
\]
The unramified character $\eta|_{G_{\mathbb Q_p}}$ is retained here;
the split-prime version of the good-twist lemma is not invoked
literally.  Potential crystallinity of $V$ follows from the actual
G68 local point, since its crystalline extension over $K_p$ is also a
finite extension of $\mathbb Q_p$.

After tensoring by $\psi$, the two graded characters have inertia
$c_p^e$ and $c_p^{-e}$.  They are therefore unramified characters
times these Tate powers and are crystalline.  The resulting extension
is potentially crystalline.  Its class maps to zero in
$H^1(B_{\rm cris}\otimes\operatorname{Hom})$ after a finite extension.
Restriction on that characteristic-zero cohomology is injective:
corestriction composed with restriction is multiplication by the
finite extension degree.  Thus its class is already crystalline over
$\mathbb Q_p$.  This proves that $V\psi$ is crystalline, including
the case of nontrivial ramified $\tau$.
Equivalently, tensoring by the actual $\tau$ removes the scalar
inertial type, and the remaining factor $\psi/\tau$ is unramified.
This is exactly the period argument of AC-K29-good-twist and the
inert version's Step~1, with the unramified determinant retained.

\emph{The exact inertia at every other rational prime.}
For $r\ne3,p$, $K/\mathbb Q$ is unramified at $r$, so rational
inertia is contained in $G_K$.  The original away-inertia theorem
\texttt{pt\_inertia\_away}, followed by the actual endpoint extraction,
gives characters $\psi,\psi^{-1}$ there.  Since $\psi^2=1$, this is
the scalar character $\psi$, with finite image and zero monodromy.
Tensoring by $\psi$ removes it exactly, including wild quadratic
inertia at $2$.

The literal \texttt{GD.cop} gives $3\nmid N$, hence $\psi$ is
unramified at $3$.  Inertia over $K_3$ acts trivially on $V$.
Rational inertia therefore acts through $C_2$, and the determinant
$\eta$ forces its two characters to be $1,\eta$.  This is finite tame
inertia with zero monodromy.  Tensoring by $\psi\eta$ preserves that
unordered pair.  Its conductor exponent is exactly one.

The primitive automorphic twist retains weight and cuspidality and
has nebentypus $\eta$, since $(\psi\eta)^2=1$.
Its representation is the displayed $R_g$, by good-prime twisting
and Chebotarev.  Away-$p$ local--global compatibility and the
classical $p$-adic compatibility statement of Saito give the full
Frobenius-semisimple Weil--Deligne parameters just computed.
At $p$, crystallinity gives an unramified parameter with zero
monodromy.  The local factor is spherical.  The primitive newform
conductor is consequently exactly $3$ by classical newvector theory.
This is an actual characteristic-zero conductor calculation;
the formal source \texttt{Level} modulus is not substituted for it.

\emph{Both primitive bad-prime roots.}
The displayed residual formula follows directly by multiplying the
two original constituents by $\psi\eta c_p^{-e}$ and using $\psi^2=1$.
At $3$, the integral projector $(1+s)/2$ for nontrivial quadratic
inertia identifies the actual invariant line.  It commutes with
reduction, whose unique inertia-trivial constituent is the global
character $1$.  Its geometric Frobenius eigenvalue is therefore
congruent to one.  The primitive Euler-factor identity gives
$a_3(g)\equiv1$.

The ordinary line of $R_g$ is unramified.  The other residual inertia
character is $\overline c_p^{-2e}$, different from one by
\texttt{GD.gen}.  Saturation of the stable line in the lattice
therefore reduces to the constituent $1$ on the full decomposition
group.  Its Frobenius eigenvalue $\alpha$ is an integral unit with
$\alpha\equiv1$.  The other Hecke root is
$\eta(p)p^{2e}/\alpha$, of positive valuation.  Thus
$a_p(g)\equiv1$ and $g$ is ordinary.  No assertion that $\psi$ itself
was unramified at $p$ has been used.
\end{proof}

\begin{corollary}[AC-G117-quadratic: conditional exclusion of the quadratic original family]
\label{thm:ACG117quadratic}
Retain the actual original-point and modularity hypotheses above, and
the checked all-$133$-residue constant/measure argument in AC-G116.
There is no such original Type II point for any $\psi$ with $\psi^2=1$.
\end{corollary}
\begin{proof}
The lemma gives the primitive level-three $g$ with precisely the
residual formula and two bad-prime coefficients used in
(G116.7)--(G116.15).  The Hecke recurrences and multiplicativity then
give the same congruence of every positive Fourier coefficient with
$E_k(1,\eta)$.  This argument is independent of the original $N$ and
$\psi$ after the primitive twist.

The weight $k=2e+1$ is odd and at least three.  The condition
\texttt{GD.gen} says $268\nmid2e$, so its residue modulo $268$
belongs to $\{3,5,\ldots,267\}$.  The exact finite certificate and
integral-measure calculation (G116.16)--(G116.20) make the constant
term of that Eisenstein series a unit at $269$ for every such weight.
The same coefficient-base-change and constant-expansion argument of
(G116.21) gives a contradiction.  No new Bernoulli computation or
new invocation of a coefficientwise imprimitive twist is required.
\end{proof}

\paragraph{Source and application boundary.}
The complete native files \texttt{Counts.lean},
\texttt{DirichletPrimeToLPartMulLPowerPart.lean},
\texttt{RhoKApply.lean}, \texttt{PtInertiaAway.lean},
\texttt{PtIneL.lean}, \texttt{AdmissibleOfTypeII.lean}, and
\texttt{Newform.lean} were freshly read for this extension.
The applied full G68 local/point proof and the good-twist proof were
also reread.  The classical compatibility and newvector inputs keep
the exact primary-source scope of G116; their deep proofs are not
newly claimed.  The all-residue certificate and measure proof are
imported from G116, with its original review record.

Nothing here excludes a representation with additional auxiliary
tame type, an arbitrary ordinary input with only the same residual
trace, or a point with changed away-kernel conditions.  Such an object
need not have the characteristic-zero conductor three used above.
Nor does this exclusion prove mixed Selmer vanishing or H04.4.
The pure arithmetic results G114--G115 remain independent of an
endpoint and are unaffected.
